%% notes-tex/025-mmli-campaign/sections/2-table.tex
\section{The campaign table\secstatus{todo}}
\label{sec:table}

\begin{table}[H]
\centering
\footnotesize
\caption{Every cell of the campaign. A number is a finished seed's score; \textit{run, ep
$k$} is the running job and the last validation epoch it logged; \textit{q$k$} is position
$k$ in the queue of Table~\ref{tab:queue}; n/a marks a seed slot the cell does not use.
Seeds a, b, c are 1, 2, 3, except in the rows inherited from the first disjoint round,
where they are 42, 1, 2. $n$ is seeds scored over seeds planned, $\Delta$ the mean minus
the comparator cell's mean (Section~\ref{sec:design}), and the best mean of each block is
bold. Generated by \texttt{mmli\_campaign\_plan.py}.}
\label{tab:campaign}
%% SOURCE: experiments/025-solid-growth/scripts/mmli_campaign_plan.py (results/mmli_campaign_plan.json) -- GENERATED, do not edit

\begin{tabular}{llllrrrrrrr}
\toprule
records & genes & fit. & config & seed a & seed b & seed c & $n$ & mean & sd & $\Delta$ \\
\midrule
\multicolumn{11}{l}{\textbf{A. Random split (010's pinned split, 37,673 validation triples)}} \\
\midrule
S0 & table & no & \texttt{s0\_r\_kl\_ctrl\_013} & 0.438 & 0.436 & 0.441 & 3/3 & 0.439 & 0.002 &  \\
S0 & table & yes & \texttt{s0\_r\_kl\_fit\_014} & 0.443 & 0.438 & 0.437 & 3/3 & 0.439 & 0.003 & +0.001 \\
S3 & table & yes & \texttt{s3\_r\_kl\_fit\_031} & 0.478 & \textit{q1} & \textit{q2} & 1/3 & \textbf{0.478} &  & +0.038 \\
S0 & composite & yes & \texttt{s0\_r\_kl\_embfit\_035} & \textit{q3} & \textit{q4} & \textit{q5} & 0/3 &  &  &  \\
S3 & composite & yes & \texttt{s3\_r\_kl\_embfit\_034} & \textit{q6} & \textit{q7} & \textit{q8} & 0/3 &  &  &  \\
S4 & table & yes & \texttt{s4\_r\_kl\_fit\_039} & \textit{q14} & \textit{q24} & \textit{q25} & 0/3 &  &  &  \\
\midrule
\multicolumn{11}{l}{\textbf{B. Query-pair-disjoint split (89 of 420 query pairs held out, 37,705 validation triples)}} \\
\midrule
S0 & table & no & \texttt{s0\_q\_kl\_ctrl\_016} & 0.140 & 0.144 & 0.125 & 3/3 & 0.136 & 0.010 &  \\
S0 & random vector & no & \texttt{s0\_q\_kl\_rand\_018} & 0.151 & 0.146 & 0.158 & 3/3 & 0.152 & 0.006 & +0.015 \\
S0 & composite & no & \texttt{s0\_q\_kl\_emb\_017} & 0.215 & 0.235 & 0.203 & 3/3 & \textbf{0.218} & 0.016 & +0.081 \\
S0 & CaLM alone & no & \texttt{s0\_q\_kl\_calm\_020} & 0.204 & \textit{run, ep 17} & \textit{q11} & 1/3 & 0.204 &  & -0.014 \\
S0 & ProtT5 alone & no & \texttt{s0\_q\_kl\_prot\_021} & 0.162 & \textit{q9} & \textit{q12} & 1/3 & 0.162 &  & -0.055 \\
S0 & flanks alone & no & \texttt{s0\_q\_kl\_fudt\_032} & 0.164 & \textit{q10} & \textit{q13} & 1/3 & 0.164 &  & -0.053 \\
S0 & table & yes & \texttt{s0\_q\_kl\_fit\_038} & \textit{q15} & \textit{q16} & \textit{q17} & 0/3 &  &  &  \\
S0 & composite & yes & \texttt{s0\_q\_kl\_embfit\_027} & 0.209 & 0.237 & 0.171 & 3/3 & 0.206 & 0.033 & -0.012 \\
S3 strict & table & yes & \texttt{s3\_q\_kl\_fit\_036} & \textit{q19} & \textit{q21} & \textit{q23} & 0/3 &  &  &  \\
S3 strict & composite & yes & \texttt{s3\_q\_kl\_embfit\_037} & \textit{q18} & \textit{q20} & \textit{q22} & 0/3 &  &  &  \\
\bottomrule
\end{tabular}

\end{table}

\notesec{What is measured}
On R, fitness supervision alone adds nothing to the triples: 0.439 with it, 0.439 without.
On Q the composite (0.218) separates from the random vector (0.152) and from the table
(0.136) with no overlap between their seeds, so sequence content carries to unseen query
pairs and a learned table does not.

\notesec{What is one seed}
The S3 seed on R scored 0.478 over epochs 10 to 29 (run yb4gjh51, job 2409033, 104 epochs).
Its validation Pearson peaked at 0.498 at epoch 32 and drifted to 0.478 by epoch 100 while
train Pearson reached 0.768, which is why every S3 and S4 cell now stops at 50 epochs.
Its row carries $n = 1/3$, so the mean and $\Delta$ printed there are that seed alone.
CaLM alone (0.204), ProtT5 alone (0.162) and the flanks alone (0.164) are one seed each.
